\documentclass[11pt,a4paper]{article}

\input{../../sablony/pripravy/priprava.sty}

\nastavPripravu{4. ročník}{Kombinatorika}{32}{Permutace I -- základy}

\begin{document}

\blok{Motivace a návaznost}

Sedm studentů si sedá do jedné řady. Kolik různých zasedacích pořádků je možných? Co se změní, když někdo musí sedět první?

\poznamka{Předpoklady: faktoriál, pravidlo součinu. Cíl: odvodit počet permutací bez opakování a řešit jednoduché úlohy s pevnou pozicí.}

\blok{Výklad a odvození}

\begin{definiceBox}
\textbf{Permutace bez opakování} je uspořádání všech $n$ navzájem různých prvků. Počet permutací značíme $P(n)$ a pro $n\in\mathbb N_0$ platí $P(n)=n!$.
\end{definiceBox}
Odvození: první pozici můžeme obsadit $n$ způsoby, druhou $n-1$ a tak dále až po jednu zbývající možnost. Při pevně určené pozici zbývají pro ostatní osoby volná místa. Zásadní je, že \textbf{řadíme všechny různé prvky}.

\blok{Společně řešené příklady}

\begin{prikladReseniBox}[title={Příklad 1}]
Kolika způsoby se do řady postaví šest různých studentů?
\tcblower
$6!=720$.
\end{prikladReseniBox}



\begin{prikladReseniBox}[title={Příklad 2}]
Pět různých knih seřadíme na polici. Kolik je pořadí, má-li učebnice matematiky být první?
\tcblower
První místo je pevně dané, zbývající čtyři knihy: $4!=24$.
\end{prikladReseniBox}



\begin{prikladReseniBox}[title={Příklad 3}]
Kolika způsoby lze usadit šest různých lidí do šesti očíslovaných židlí, jestliže Eva sedí na poslední židli?
\tcblower
Pět zbývajících obsadí pět zbylých míst: $5!=120$.
\end{prikladReseniBox}



\blok{Samostatná práce}

\begin{zadaniReseniBox}[title={Samostatná práce 1}]
Seřaď čtyři různá písmena A, B, C, D. Kolik pořadí existuje?
\tcblower
$4!=24$.
\end{zadaniReseniBox}



\begin{zadaniReseniBox}[title={Samostatná práce 2}]
Kolika způsoby lze postavit sedm lidí do fronty, musí-li Petr stát první?
\tcblower
$6!=720$.
\end{zadaniReseniBox}



\begin{zadaniReseniBox}[title={Samostatná práce 3}]
Kolik pořadí vznikne ze šesti různých předmětů, musí-li modrý předmět stát poslední?
\tcblower
$5!=120$.
\end{zadaniReseniBox}



\blok{Rozšíření}

\begin{zadaniReseniBox}[title={Rozšíření 1}]
V řadě sedmi lidí musí být Anna první a Boris poslední. Kolik pořadí je možných?
\tcblower
Obě místa jsou pevná, zůstává pět lidí: $5!=120$.
\end{zadaniReseniBox}



\blok{Domácí úkol}

\poznamka{Doporučená práce doma: přibližně 15--30 minut. Úlohy jsou seřazeny od základních k náročnějším; učitel může poslední podúlohu označit jako volitelnou.}

\begin{zadaniReseniBox}[title={Domácí úkol}]

\begin{enumerate}[label=\alph*),leftmargin=*,itemsep=2pt]

\item Vypočti počet pořadí pěti různých knih.
\item Kolika způsoby se seřadí osm různých osob, pokud je konkrétní osoba první?
\item V soutěži vystupuje šest různých týmů. Kolik pořadí vystoupení existuje, pokud tým A vystupuje poslední?
\item Kolik pořadí sedmi různých předmětů vznikne, pokud jeden označený předmět musí být první a jiný poslední?
\item Osm různých studentů stojí v řadě. Anna musí stát první a Boris druhý. Kolik je pořadí?
\item Ze sedmi různých lidí musejí Eva a Marek stát na prvních dvou místech v libovolném vzájemném pořadí. Kolik výsledků?

\end{enumerate}

\tcblower

\begin{enumerate}[label=\alph*),leftmargin=*,itemsep=2pt]

\item $5!=120$.
\item $7!=5040$.
\item $5!=120$.
\item $5!=120$.
\item $6!=720$.
\item $2!\cdot5!=240$.

\end{enumerate}

\end{zadaniReseniBox}

\poznamka{Orientační práce v hodině: 0--12 min odvození; 12--27 min ukázky; 27--42 min samostatná práce; 42--45 min kontrola.}

\poznamka{Obtížnější bloková metoda zůstává až do následující hodiny; nepřidávej ji před zvládnutím základní situace.}

\end{document}
